\def\stat{spivak}



\def\tit{МЕТАОБЪЕКТНЫЙ ПОДХОД К~МОДЕЛИРОВАНИЮ БИЗНЕС-ПРОЦЕССОВ 

ПРЕДПРИЯТИЯ\\ В~РАМКАХ ЕДИНОЙ ERP-СИСТЕМЫ}



\def\titkol{Метаобъектный подход к~моделированию бизнес-процессов 

предприятия} % в~рамках единой ERP-системы}





\def\aut{С.\,И.~Спивак$^1$, Н.\,Д.~Морозкин$^2$, Л.\,А.~Лукьянов$^3$}



\def\autkol{С.\,И.~Спивак, Н.\,Д.~Морозкин, Л.\,А.~Лукьянов}



\titel{\tit}{\aut}{\autkol}{\titkol}



\index{Спивак С.\,И.}

\index{Морозкин Н.\,Д.}

\index{Лукьянов Л.\,А.}

\index{Spivak S.\,I.}

\index{Lukyanov L.\,A.}

\index{Morozkin N.\,D.}





%{\renewcommand{\thefootnote}{\fnsymbol{footnote}}

%\footnotetext[1]

%{Работа выполнена при поддержке гранта Президента РФ (МД-707.2017.8) 

%и~Госзадания (соглашение  №\,2.3440.2017/4.6).}}



\renewcommand{\thefootnote}{\arabic{footnote}}

\footnotetext[1]{Башкирский государственный университет, semen.spivak@mail.ru}

\footnotetext[2]{Башкирский государственный университет, rector@bsunet.ru}

\footnotetext[3]{Башкирский государственный университет, leonlukyanov@gmail.com}





 \label{st\stat}



                  \thispagestyle{headings}

                  

\vspace*{-6pt}

  

  

  

  

\Abst{Объединение бизнес-процессов предприятия внутри единой информационной системы 

носит массовый характер. Оптимизация дея\-тель\-ности предприятий в~условиях гибкого 

рынка и~стремительного технического прогресса невозможна без использования единой 

информационной системы. Существует огромное число ERP

(Enterprise Resource Planning) сис\-тем, призванных решить 

в~полной мере задачу автоматизации всех процессов на предприятии. Однако ни одна из них 

не способна в~полной мере учитывать существующие механизмы автоматизации на 

предприятии. Эти механизмы обычно реализованы не только в~виде документов 

с~техническим заданием к~тому или иному бизнес-процессу, но и~как программные решения, 

обычно разработанные с~помощью достаточно устаревших технологий. В~данной работе 

рассматривается ERP-сис\-те\-ма собственной разработки, в~которую заложен абсолютно 

новый подход внедрения, разработки и~организации единого информационного пространства 

для всех подразделений предприятия. Предлагается метод внедрения, основанный на 

последовательном замещении существующей реализации биз\-нес-про\-цес\-сов подсистемами 

внедряемой платформы. Рассматривается метаобъектный подход к~реализации  

биз\-нес-про\-цес\-сов, который может быть легко использован специалистами предметных 

областей, работающими непосредственно на предприятии. Предлагаемая ERP-сис\-те\-ма не 

изменяет логику существующих биз\-нес-про\-цес\-сов, но обеспечивает их интеграцию, 

прозрачность и~единую модель данных.}



\KW{MRP; MRP II; MES; CRUD; единое информационное пространство; планирование 

ресурсов; моделирование информационных систем; платформа; архитектура ERP}



 \DOI{10.14357\!/\!08696527190210} 





\vskip 10pt plus 9pt minus 6pt





      \thispagestyle{headings}

      

      \vspace*{-9pt} 



\section{Введение}



%\vspace*{-2pt}



  Как известно, существует огромное число информационных систем, 

призванных регулировать деятельность подразделений на предприятиях 

различной отрас\-ле\-вой принадлежности: промышленной, торговой, 

транспортной, финан-\linebreak\vspace*{-12pt}



\pagebreak



\noindent

совой. Такие информационные системы помимо основной 

функциональности оптимизируют деятельность того или иного подразделения 

на предприятии. Очевидно, что отделы в~организации постоянно 

взаимодействуют между собой в~процессе создания товаров и~услуг при 

реализации заказа. Именно по этой причине необходима интеграция 

информационных систем (особенно на крупном предприятии) для обеспечения 

прозрачности всех бизнес-процессов и~оптимизации деятельности всей 

организации в~целом~[1]. 

  

  В настоящее время понятие интегрированных подсистем заложено 

в~маркетинговом термине <<ERP-сис\-те\-ма>>. Широко распространено 

следующее определение этого термина: ERP~--- <<это организационная 

стратегия интеграции производства и~операций, управления трудовыми 

ресурсами, финансового менеджмента и~управления активами, 

ориентированная на непрерывную балансировку и~оптимизацию ресурсов 

предприятия посредством специализированного интегрированного пакета 

прикладного программного обеспечения, обес\-пе\-чи\-ва\-юще\-го общую модель 

данных и~процессов для всех сфер деятельности>>~[2]. С~общей точки зрения 

под ERP-сис\-те\-мой понимают интегрированный пакет программного 

обеспечения, включающий в~себя совокупность модулей (под\-сис\-тем), задача 

которых состоит в~математическом и~информационном обеспечении 

определенного биз\-нес-про\-цес\-са на предприятии. Исторически сложилось, 

что ERP-сис\-те\-мы успешно решают задачу интеграции и~оптимизации всех 

биз\-нес-про\-цес\-сов организации. Первые ERP-сис\-те\-мы представляли собой 

системы планирования материалов. На сегодняшний день уже разработаны 

ERP-сис\-те\-мы третьего поколения, организующие веб-сер\-ви\-сы для 

мобильных устройств~[2]. Например, такая подсистема может 

автоматизировать работу отдела маркетинга, финансового отдела или отдела 

ма\-те\-ри\-аль\-но-тех\-ни\-че\-ско\-го обеспечения, особенно на предприятиях 

производственного типа. Цель ERP-сис\-те\-мы состоит в~автоматизации всех 

информационных процессов как внутри самой организации, так и~внутри ее 

контрагентов. К~недостаткам современных ERP-сис\-тем можно от\-нес\-ти 

сложность их внедрения и~высокие финансовые затраты как при покупке, так 

и~в процессе работы с~данной системой~[3]. А~к~достоинствам ERP-сис\-тем 

следует отнести прозрачность всех биз\-нес-про\-цес\-сов, простоту 

в~управлении и~устранение человеческого фактора. 

  

  Стоит отметить, что обычно при внедрении ERP-сис\-тем не производится 

интеграция с~существующей системой; таким образом, весь накопленный опыт 

предприятия используется только на этапе внедрения новой системы 

посредством словесного или документального описания сотрудниками 

особенностей сложившихся на предприятии биз\-нес-про\-цес\-сов 

поставщикам данной ERP-сис\-те\-мы. Можно наблюдать ситуации, когда  

биз\-нес-де\-я\-тель\-ность подстраивается под ERP-сис\-те\-му. Такой подход 

можно считать вполне адекватным и~логичным в~случае, когда организация 

небольшая или находится в~процессе становления. Однако для крупных 

предприятий важно сохранить сложившиеся особенности деятельности всех 

подразделений, с~тем чтобы не потерять накопленный опыт в~вопросах 

управления, планирования и~кооперации. 

  

  Авторы предлагают концептуально новый подход к~моделированию  

биз\-нес-про\-цес\-сов любого предприятия с~использованием инструмента 

собственной разработки~--- MOPEX (Meta Object Process EXecution). 

Моделирование включает как описательную, так и~практическую части. 

Инструмент имеет гибкую архитектуру, множество вспомогательных сервисов 

и~огромный потенциал к~расширению и~оптимизации ядра новыми 

возможностями. Математическая составляющая сервисов для решения задач 

планирования и~управления пред\-став\-ле\-на в~виде подключаемых модулей 

различной структуры. Главная особенность разработанной платформы 

заключается в~простоте моделирования новых, а~также переработки 

существующих информационных процессов для автоматизации деятельности 

подразделений предприятия. 





\begin{figure}[b] %fig1

  \vspace*{-12pt}

 \begin{center}

 \mbox{%

 \epsfxsize=127.909mm 

 \epsfbox{spi-1.eps}

 }

 \end{center}

\vspace*{-15pt}

\Caption{Архитектура платформы MOPEX}

\end{figure}



\vspace*{-3pt}



\section{Архитектура MOPEX}



\vspace*{-4pt}



  Архитектура платформы MOPEX, представленная на рис.~1, имеет 

достаточно сложную структуру, обусловленную требованиями гибкости 

разработки, безопасности и~общей производительности. Взаимодействие 

основных компонентов архитектуры~--- сервера приложений, клиента 

пользователя, клиента администратора~--- осуществляется по сети согласно 

трехзвенной схеме~[4]. Данная платформа позиционируется в~качестве 

высокопроизводительной расширяемой ERP-сис\-те\-мы, которая 

предоставляет инструмент описания биз\-нес-про\-цес\-сов. Стандартные 

решения, включаемые большинством ERP-сис\-тем, предлагаемых на рынке  

ИТ-ин\-дуст\-рии, в~данной платформе отсутствуют. Связано это с~тем фактом, 

что предприятие при внедрении данной платформы должно самостоятельно 

моделировать биз\-нес-про\-цес\-сы по метаобъектной модели разработки.

  



  Особенность архитектуры платформы заключается в~наличии стека плагинов 

по управлению биз\-нес-ло\-ги\-кой и~стека расширений для графической 

обработки информации на стороне пользовательских приложений. Плагины, 

как и~расширения, независимы друг от друга и~могут подключаться 

в~зависимости от роли пользователя системы. Например, если пользователь 

является разработчиком предметной области, то ему должен быть доступен 

инструмент построения отчетов и,~возможно, инструмент создания 

пользовательского интерфейса. В~противоположность этому, пользователь, 

выполняющий роль ответственного за ведение задачи, должен иметь доступ 

только к~инструменту отображения отчетов. Набор функциональных 

возможностей строго определяется ролями пользователей и~описывается 

дополнительными атрибутами (разрешен доступ, запись, выполнение, др.), 

позволяющими тонко настроить способ доступа к~данным.

  

\section{Метаобъектная модель разработки}



  Прежде чем перейти к~рассмотрению архитектуры платформы, необходимо 

описать процесс разработки подсистемы для решения некоторой задачи 

предметной области. Следует отличать модуль ERP-системы от 

<<подсистемы>> предметной области. Модуль ERP-сис\-те\-мы~--- это 

автономный программный продукт, предлагающий стандартное решение для 

автоматизации работы отдельного подразделения предприятия. Например, 

модулем ERP-системы является программа внутрицехового планирования~[5]. 

В~противоположность модулю, подсистема предметной области~--- это 

программный продукт, автоматизирующий некоторый биз\-нес-про\-цесс. 

К~подсистемам, например, можно отнести <<Ведение маршрутов деталей 

и~сборочных единиц (ДСЕ) внутри цеха>>, <<Управление ДСЕ>>.

  

  Для реализации подсистем на платформе MOPEX предлагается 

метаобъектная модель разработки, представленная на рис.~2. Данная модель 

представляет собой модификацию широко известной итеративной модели 

разработки~[6]. Метаобъектная модель разработки осуществляется 

специалистами предметной области и~включает в~себя следующие этапы. 

\begin{enumerate}[1.]

\item \textit{Разработка информационной модели}. Формируется 

техническое задание (ТЗ) на разрабатываемую подсистему сотрудниками 

отдела. Осуществляется работа по описанию структуры данных 

разрабатываемой подсистемы. Структура данных~--- это набор связанных 

таблиц, который регламентирует способ хранения данных подсистемы. 

Необходимо отметить, что на этапе создания информационной модели 

следует учитывать, что одни и~те же данные могут быть использованы 

в~разных подсистемах. Например, при создании подсистемы <<Ведение 

маршрутов ДСЕ внут\-ри цеха>> структура данных 

включает в~себя данные по деталям, операциям и~оборудованию.



\begin{figure} %fig2

\vspace*{1pt}

 \begin{center}

 \mbox{%

 \epsfxsize=93.871mm 

 \epsfbox{spi-2.eps}

 }

 \end{center}

\vspace*{-9pt}

  \Caption{Метаобъектная модель разработки}

  \end{figure}

  

  

\item \textit{Разработка функциональной модели}. Определяются функции 

разрабатываемой подсистемы, строятся диаграммы USE CASE, IDEF0, 

DFD. Например, в~подсистеме <<Ведение маршрутов ДСЕ 

внут\-ри цеха>> можно выделить следующие функции: просмотр 

справочника операций цеха, создание маршрута движения ДСЕ, 

корректировка загруженности оборудования и~т.\,д.

\item \textit{Реализация информационной модели}. Информационная модель 

создается в~системе MOPEX в~виде метаобъектов. Создание осуществляется 

с~использованием клиента <<Администратор>> платформы MOPEX. 

Метаобъект~--- это класс и~его окружение. Классом может выступать любой 

объект подсистемы, наделенный атрибутами. Окружение класса 

характеризует связи с~другими классами подсистемы и~иногда определяет 

встроенные функции, выполняемые при взаимодействии с~данным классом. 

На рис.~3 приведен пример класса ДСЕ, используемого в~подсистеме 

<<Управление ДСЕ>>.



\begin{figure} %fig3

  \vspace*{1pt}

 \begin{center}

 \mbox{%

 \epsfxsize=128mm 

 \epsfbox{spi-3.eps}

 }

 \end{center}

\vspace*{-9pt}

  \Caption{Пример класса информационной модели}

  \end{figure}

  





\item \textit{Реализация функциональной модели}. На основе 

функциональной модели производится описание взаимодействия 

метаобъектов средствами платформы MOPEX, а именно выполняются 

следующие последовательные шаги:

\begin{itemize}

\item[(а)] создание ролей подсистемы в~соответствии с~диаграммой USE 

CASE;

\item[(б)]создание пользователей подсистемы;

\item[(в)]назначение пользователям определенных ролей;

\item[(г)]создание пунктов и~подпунктов меню подсистемы, отражающих 

основные функции подсистемы;

\item[(д)]определение контекстного меню, обеспечивающего выполнение 

некоторой функции, зависящей от текущего набора данных. 

\end{itemize}



 На рис.~4 приведен пример реализации функциональной модели 

подсистемы <<Управление ДСЕ>>. 

\item \textit{Создание окружения подсистемы}. Производится создание 

дополнительных элементов разрабатываемой подсистемы: отчетных форм, 

экранных форм, функций по умолчанию, пакетных задач. Стоит отметить, 

что данные элементы не всегда являются обязательными в~той или иной 

подсистеме. Например, в~подсистеме <<Управление ДСЕ>> отчетные 

и~экранные формы отсутствуют.

\item \textit{Апробация подсистемы}. Осуществляется тестирование 

подсистемы специалистами подразделения, для которого 

разрабатывалась данная подсистема. При выявлении недоработок 

подсистемы производится уточнение ошибок в~форме дополнения 

к~основному ТЗ. Затем проводится доработка подсистемы, начиная 

с~первого этапа метаобъектной модели разработки.

\end{enumerate}



  \begin{figure} %fig4

  \vspace*{1pt}

 \begin{center}

 \mbox{%

 \epsfxsize=128mm 

 \epsfbox{spi-4.eps}

 }

 \end{center}

\vspace*{-9pt}

  \Caption{Пример описания функциональной модели}

  \end{figure}





  



\section{Заключение}



  В работе были исследованы существующие на рынке платформы класса ERP; 

выявлены недостатки предлагаемых решений; предложена актуальная 

архитектура платформы для построения биз\-нес-при\-ло\-же\-ний; предложена 

модель разработки, покрывающая все этапы создания бизнес-приложений~--- от 

проектирования до внедрения в~эксплуатацию. Актуальность и~новизна 

платформы среди всех существующих решений подкрепляется методом 

представления объекта информации и~его контекста взаимодействия в~виде 

метаобъекта, который при необходимости может быть использован простым 

внедрением в~другом биз\-нес-про\-цессе без копирования или преобразования 

со стороны разработчика задачи предметной области.

  

{\small\frenchspacing

{%\baselineskip=10.8pt

\addcontentsline{toc}{section}{Литература}

\begin{thebibliography}{9}

\bibitem{1-spi}

\Au{Григорьев А.\,А., Титов В.\,А.} Характеристика, структура, организация систем 

управ\-ле\-ния ERP, ERP~II и~ERP~III~/\!\!/ Фундаментальные исследования, 2017. №\,2. 

С.~48--51.

\bibitem{2-spi}

\Au{Abdullah A.\,M.\,A.} Evolution of enterprise resource planning~/\!\!/ 

Excel J.~Eng. Technology Management Sci., 2017. Vol.~1. No.\,11. P.~1--6.

\bibitem{3-spi}

\Au{Елисеев В.\,Г., Фисенко~Н.\,И.} Анализ затрат на внедрение ERP систем~/\!\!/ Научная 

сессия МИФИ-2007: Сборник научных трудов.~--- М.: МИФИ, 2007. Т.~2. С.~48--49.

\bibitem{4-spi}

\Au{Шмидт И.\,А.} Архитектура платформы для разработки биз\-нес-при\-ло\-же\-ний~/\!\!/ 

Современные проблемы науки и~образования, 2014. №\,6. С.~348.

\bibitem{5-spi}

\Au{Васильева Ю.\,В.} Планирование производственных ресурсов (MRP~II) и~планирование 

потребностей предприятия (ERP)~/\!\!/ Современные тенденции в~экономике 

и~управ\-ле\-нии: новый взгляд.~--- Новосибирск: Центр развития научного сотрудничества, 

2015. С.~46--50.

\bibitem{6-spi}

\Au{Терехов А.\,Н.} Технология программирования.~--- М.: НОУ ИНТУИТ, 2006. 77~с.

        \end{thebibliography}

} }





\vspace*{-3pt}



\hfill{\small\textit{Поступила в~редакцию 02.05.18}}









\vspace*{8pt}



\hrule



\vspace*{2pt}



%\newpage



%\vspace*{-24pt}



\hrule



%\vspace*{9pt}







\def\tit{METAOBJECT APPROACH FOR~MODELING ENTERPRISE BUSINESS PROCESSES INSIDE 

SINGLE ERP SYSTEM}



\def\titkol{Metaobject approach for~modeling enterprise business processes inside 

single ERP system}





\def\aut{S.\,I.~Spivak, L.\,A.~Lukyanov, and~N.\,D.~Morozkin}

\def\autkol{S.\,I.~Spivak, L.\,A.~Lukyanov, and~N.\,D.~Morozkin}







\titel{\tit}{\aut}{\autkol}{\titkol}



\vspace*{-9pt}



\noindent

\begin{center}

Bashkir State University, 32~Validy Str., Ufa 450076, Russian Federation

\end{center}





 \noindent



\def\leftfootline{\small{\textbf{\thepage}

\hfill Sistemy i Sredstva Informatiki~---

Systems and Means of Informatics 2019 vol~29 no\,2}

}%

 \def\rightfootline{\small{Sistemy i Sredstva Informatiki~---

 Systems and Means of Informatics   2019   vol~29   no\,2

\hfill \textbf{\thepage}}}



   

\Abste{Nowadays, combining enterprise business processes within a~single information system has 

a~widespread trend. Work optimization of enterprises in condition of flexible market and rapid 

technical progress is impossible without using a~single information system.

 There are many ERP (Enterprise Resource Planning)

systems designed to solve the main problem of automation of all enterprise business processes. 

However, neither one of them is able to take into account existing automation mechanisms of the 

enterprise. These mechanisms are often designed not only using documents with technical order for 

any type of business process but also using software developed with outdated technics. In this 

paper, the authors consider their own ERP system, which contains a~completely new approach to 

the application, development, and organization of a~single information space for all business units. 

The proposed approach is based on sequential replacement of the existing implementation of 

business processes by platform subsystems. The metaobject approach to the implementation of 

business processes is considered, which is applied by specialists of subject areas working directly in 

the enterprise. The proposed ERP system does not change the logic of existing business processes 

but ensures their integration, transparency, and common data model.}

 

\KWE{MRP; MRPII; ERP; MES; CRUD; common information space; resource planning; 

information system modeling; platform; software architecture}



 \DOI{10.14357\!/\!08696527190210} 



%\vspace*{-24pt}



%\Ack

%\noindent

 





\renewcommand{\bibname}{\protect\rmfamily References}

%\renewcommand{\bibname}{\large\protect\rm References}



{\small\frenchspacing

{%\baselineskip=10.8pt

\addcontentsline{toc}{section}{References}

\begin{thebibliography}{9}

\bibitem{1-spi-1}

\Aue{Grigorev, A.\,A., and V.\,A.~Titov.} 

2017. Kharakteristika, struktura, organizatsiya 

sistem upravleniya ERP, ERP~II i~ERP~III [Characteristics, structure, 

organization systems ERP-management, ERP~II and ERP~III]. \textit{Fundamental Research} 2:48--51.

\bibitem{2-spi-1}

\Aue{Abdullah, A.\,M.\,A.} 2017. Evolution of enterprise resource 

planning. \textit{Excel J.~Eng. Technology Management Sci.} 

1(11):1--6.

\bibitem{3-spi-1}

\Aue{Eliseev, V.\,G., and N.\,I.~Fisenko}.

 2017. Analiz zatrat na vnedrenie ERP sistem 

[Сost analysis for the introduction of ERP systems]. 

\textit{Nauchnaya sessiya MIFI-2007: Sbornik nauchnykh trudov} 

[Scientific session MEPhI-2007: Collection of scientific papers]. 2:48--49.

\bibitem{4-spi-1}

\Aue{Shmidt, I.\,A.} 2014. Arkhitektura platformy dlya razrabotki biznes-prilozheniy 

[The architecture of a~development platform for business applications]. 

\textit{Modern Problems of Science 

and Education} 6:348.

\bibitem{5-spi-1}

\Aue{Vasil'eva, Yu.\,V.} 2015. Planirovanie proizvodstvennykh resursov (MRP~II) 

i~pla\-ni\-ro\-va\-nie potrebnostey predpriyatiya (ERP) [Resource planning (MRP~II) 

and enterprise needs planning (ERP)]. \textit{Sovremennye tendentsii v~ekonomike 

i~upravlenii: novyy vzglyad} [Current trends in economics and management: 

A~modern point]. Novosibirsk: Center for the Development of Scientific Cooperation. 46--50.

\bibitem{6-spi-1}

\Aue{Terekhov, A.\,N.} 2006. \textit{Tekhnologiya programmirovaniya} [Programming 

technology]. Moscow: NOU INTUIT. 77~p.

\end{thebibliography}

} }





\vspace*{-3pt}



\hfill{\small\textit{Received May 2, 2018}} 



\Contr



\noindent

\textbf{Spivak Semen I.} (b.\ 1945)~--- Doctor of Science in physics and 

mathematics, professor, Head of Department, Bashkir State University, 32~Validy 

Str., Ufa 450076, Russian Federation; \mbox{semen.spivak@mail.ru}



\vspace*{3pt}





\noindent

\textbf{Morozkin Nikolay D.} (b.\ 1953)~--- Doctor of Science in physics and 

mathematics, Chancellor, Bashkir State University, 32~Validy Str., Ufa 450076, 

Russian Federation; \mbox{rector@bsunet.ru}



\vspace*{3pt}



\noindent

\textbf{Lukyanov Leonid A.} (b.\ 1992)~--- PhD student, Bashkir State 

University, 32~Validy Str., Ufa 450076, Russian Federation; 

\mbox{leonlukyanov@gmail.com}

\label{end\stat}





\renewcommand{\bibname}{\protect\rm Литература}

